%=====================================================================
% The Webb Dipole and Spatial Variation of Constants
% Converted from standalone paper: webb_dipole.tex
%=====================================================================
\section{The Webb Dipole and Spatial Variation of Constants}
\label{ch:webb_dipole}

%=====================================================================
\subsection{Introduction: Feynman's Question}
%=====================================================================

Richard Feynman called the fine structure constant $\alpha_\text{em} \approx 1/137.036$ a ``magic number'' that all physicists should lose sleep over. For 70 years, the question ``why 1/137?'' has had no answer. The Standard Model treats it as an input parameter --- a number that must be measured, not derived.

In 1999, Webb et al.\ reported evidence that $\alpha_\text{em}$ is not even constant: quasar absorption spectra suggest a spatial variation $\Delta\alpha/\alpha \sim 10^{-5}$ with a dipole pattern across the sky \cite{Webb}. This finding, if confirmed, would overturn the assumption that physical laws are identical everywhere. For 25 years, the community has been split: is the Webb dipole real physics or systematic error?

We propose a resolution from Dynamic Resolution Lattice Theory (DRLT) : \textbf{$\alpha_\text{em}$ is not a fundamental constant. It is a local manifestation of $\det(G_h)$ redistribution.} The true constant is $\agut = 6/(25\pi^2)$, derived from lattice geometry. The measured $\alpha_\text{em}$ varies because the geometric contributions $\Delta_i$ --- $\det(G_h)$ weights from the rank projection --- are spatially non-uniform. Both sides of the Webb debate are correct: $\alpha_\text{em}$ varies (Webb is right), but the underlying physics is universal ($\agut$ is constant).

%=====================================================================
\subsection{Trace Conservation (Review)}
%=====================================================================

\subsubsection{Coupling constants from DRLT}

In DRLT , the three gauge couplings at $M_Z$ are derived from the Binet--Cauchy decomposition of $\Lambda^3(\CC^5)$:
\begin{align}
1/\alpha_3 &= 1 \times 8 \times S(1) = 8.0, \\
1/\alpha_2 &= 12 \times 2 \times S(2) = 30.0, \\
1/\alpha_1 &= 12 \times 3 \times S(\infty) = 6\pi^2 \approx 59.2,
\end{align}
where $S(N) = \sum_{n=1}^{N} 1/n^2$ is the partial Basel sum.

These combinatorial values receive $\det(G_h)$ geometric contributions:
\begin{align}
\delta_3 &= (1/\alpha_3)_\text{obs} - 8.0 = +0.47, \\
\delta_2 &= (1/\alpha_2)_\text{obs} - 30.0 = -0.40, \\
\delta_1 &= (1/\alpha_1)_\text{obs} - 59.2 = -0.22.
\end{align}

\subsubsection{The sum rule}

The geometric contributions obey trace conservation:
\begin{equation}
\sum_i \delta_i = \delta_3 + \delta_2 + \delta_1 + \delta_G = +0.47 - 0.40 - 0.22 + 0.15 = 0.00,
\end{equation}
proven from trace conservation: $\tr(G) = N$ is invariant under the unitary projection from rank $N$ to rank 5.

\subsubsection{The base perturbation $\varepsilon_0$}

All geometric contributions scale with a single parameter $\varepsilon_0 \approx 0.0038$:
\begin{equation}
\delta_i = \text{Sgn}_i \times (1/\alpha_i)_\text{pred} \times M_i \times \varepsilon_0,
\end{equation}
where $M_i$ are geometric weights determined by $(n_S, n_T) = (3, 2)$. The parameter $\varepsilon_0$ is not free: it encodes the local ``diagonalization progress'' of the Gram matrix --- the cosmic address.

%=====================================================================
\subsection{Spatial Variation of $\det(G_h)$ Contributions}
%=====================================================================

\subsubsection{Key insight: $\varepsilon_0$ is a local field}

In the simplex network (Chapter~10), each spatial location corresponds to a set of simplices with specific $\psi$ values at their vertices. The Gram determinant $\det(G_h)$ varies from simplex to simplex because the $\psi$ values differ. The geometric parameter $\varepsilon_0$ is a function of the local $\det(G_h)$ distribution, so it inherits spatial dependence:
\begin{equation}
\varepsilon_0 = \varepsilon_0(\mathbf{x}).
\end{equation}
This is not an extrapolation: it is a direct consequence of the simplex network having different $\psi$ configurations at different locations. The block universe is a static assignment of complex numbers to vertices; different regions of the block have different assignments, hence different $\varepsilon_0$.

\subsubsection{What varies and what doesn't}

\begin{theorem}[Spatial trace conservation]
At every spatial position $\mathbf{x}$:
\begin{equation}
\sum_i \delta_i(\mathbf{x}) = 0.
\end{equation}
The total coupling $1/\agut = 25\pi^2/6$ is spatially invariant. Only the distribution of geometric weight among sectors varies.
\end{theorem}

\begin{proof}
$\tr(G) = N$ is a global constraint, independent of position. The Binet--Cauchy completeness $\sum_\text{channels} = d^2 = 25$ is a geometric identity, independent of the state $\psi$. Therefore $\sum_i(1/\alpha_i) = 25\pi^2/6$ at every point, and $\sum_i \delta_i = 0$ everywhere.
\end{proof}

\begin{corollary}
The fine structure constant
\begin{equation}
1/\alpha_\text{em} = \tfrac{5}{3}(1/\alpha_1) + (1/\alpha_2) = \tfrac{5}{3}(59.2 + \delta_1) + (30.0 + \delta_2)
\end{equation}
varies spatially because $\delta_1(\mathbf{x})$ and $\delta_2(\mathbf{x})$ vary, even though their contribution to $\agut$ is fixed.
\end{corollary}

%=====================================================================
\subsection{Quantitative Comparison with the Webb Dipole}
%=====================================================================

\subsubsection{Observed signal}

Webb et al.\ report a dipole variation in $\alpha_\text{em}$ across the sky:
\begin{equation}
\frac{\Delta\alpha}{\alpha} \approx (1.1 \pm 0.25) \times 10^{-5},
\end{equation}
with the dipole axis pointing toward $(\text{RA}, \text{Dec}) \approx (17.3^\text{h}, -61^\circ)$ \cite{Webb}.

\subsubsection{DRLT prediction}

If $\varepsilon_0$ varies by a fraction $f = \Delta\varepsilon_0/\varepsilon_0$ across the sky:
\begin{align}
\Delta\delta_1 &= \delta_1 \times f = -0.22\,f, \\
\Delta\delta_2 &= \delta_2 \times f = -0.40\,f, \\
\Delta(1/\alpha_\text{em}) &= \tfrac{5}{3}\,\Delta\delta_1 + \Delta\delta_2 = -0.767\,f.
\end{align}

The fractional variation:
\begin{equation}
\frac{\Delta\alpha_\text{em}}{\alpha_\text{em}} = \frac{0.767\,f}{128.7} = 5.96 \times 10^{-3}\,f.
\end{equation}

Setting this equal to the observed $10^{-5}$:
\begin{equation}
\boxed{f = \frac{10^{-5}}{5.96 \times 10^{-3}} = 1.68 \times 10^{-3} = 0.17\%.}
\end{equation}

\subsubsection{Comparison with the CMB dipole}

The CMB dipole (our velocity relative to the CMB rest frame) gives:
\begin{equation}
\frac{v}{c} = \frac{370\;\text{km/s}}{3 \times 10^5\;\text{km/s}} = 1.23 \times 10^{-3} = 0.12\%.
\end{equation}

The required $\varepsilon_0$ variation (0.17\%) is \textbf{1.4 times the CMB dipole} --- the same order of magnitude. This is non-trivial: the geometric parameter varies at a scale consistent with the known large-scale anisotropy of the universe.

\subsubsection{Direction}

The Webb dipole axis $(17.3^\text{h}, -61^\circ)$ does not align with the CMB dipole $(11.2^\text{h}, -7^\circ$; $\sim 100^\circ$ separation). In DRLT, this is expected: $\varepsilon_0$ depends on the local $\det(G_h)$ distribution (large-scale structure), which has its own dipole independent of our peculiar velocity.

%=====================================================================
\subsection{Falsifiable Predictions}
%=====================================================================

The $\det(G_h)$ redistribution mechanism makes four testable predictions that distinguish it from all other models of varying constants:

\subsubsection{Prediction 1: $\alpha_s$ anti-correlates with $\alpha_\text{em}$}

The spatial trace conservation $\sum_i \Delta_i(\mathbf{x}) = 0$ requires that if $\alpha_\text{em}$ is larger in some direction, $\alpha_s$ must be \textbf{smaller} in that direction:
\begin{equation}
\Delta\delta_3 = -(\Delta\delta_1 + \Delta\delta_2 + \Delta\delta_G) \approx -\Delta\delta_1 - \Delta\delta_2.
\end{equation}

This is testable by directly measuring $\alpha_s$ variations in quasar spectra.

\begin{remark}[$\mu = m_p/m_e$ is trace-protected]
An earlier version of this analysis predicted $\Delta\mu/\mu$ would anti-correlate with $\Delta\alpha_\text{em}/\alpha_\text{em}$ through the chain $\alpha_s \to \Lambda_\text{QCD} \to m_p$. However, DRLT's own results (Chapter 6) show this chain is broken: $\Lambda_\text{QCD} = v_H/\sqrt{c} \times \agut^2 \times n_A$ depends on $\agut$ (trace-invariant), not on $\alpha_s$ (which varies with $\varepsilon_0$). Similarly, $m_e$ depends on the lattice constant $\delta_\text{EW} = 1 - S(2)/S(\infty)$. Therefore $\mu \approx \text{const}$ across the sky: $\Delta\mu/\mu \approx 0$. Published H$_2$ quasar data ($\sim 10$ systems, all within $1$--$2\sigma$ of zero) are consistent with this prediction.
\end{remark}

\subsubsection{Prediction 2: $\agut$ is spatially invariant}

The combination
\begin{equation}
\frac{1}{\agut} = \frac{1}{\alpha_3} + \frac{12 \times 2}{1 \times 8}\frac{1}{\alpha_2} + \frac{12 \times 3}{1 \times 8}\frac{1}{\alpha_1} = \frac{25\pi^2}{6}
\end{equation}
must be identical at every spatial location. Any observed variation in this combination would falsify the trace conservation mechanism.

\subsubsection{Prediction 3: Dipole amplitude depends on redshift}

The geometric parameter $\varepsilon_0(z)$ changes with cosmic epoch (redshift $z$). The dipole amplitude should therefore vary with $z$:
\begin{equation}
\frac{\Delta\alpha}{\alpha}(z) = 5.96 \times 10^{-3} \times \frac{\Delta\varepsilon_0(z)}{\varepsilon_0(z)}.
\end{equation}

At high $z$ (early universe, larger $\varepsilon_0$), the fractional variation $\Delta\varepsilon_0/\varepsilon_0$ could differ, producing a $z$-dependent dipole. This is testable with high-$z$ quasar surveys.

\subsubsection{Prediction 4: No variation in $\alpha_\text{em}$ at the GUT scale}

At energies approaching $M_\text{GUT}$, all forces see $N_\text{eff} = \infty$, geometric weight is uniformly distributed, and $\varepsilon_0 \to 0$. Therefore $\Delta\alpha/\alpha \to 0$ at high energy. If constant variation is observed in early-universe processes (e.g., BBN), DRLT is falsified.

%=====================================================================
\subsection{The True Constant}
%=====================================================================

Feynman asked: ``Why 1/137?'' DRLT answers: \textbf{the question is wrong.}

$1/137$ is not a fundamental constant. It is a local value of $1/\alpha_\text{em}$, determined by the $\det(G_h)$ redistribution at our cosmic address ($\varepsilon_0 \approx 0.0038$). The fundamental constant is:
\begin{equation}
\boxed{\frac{1}{\agut} = \frac{25\pi^2}{6} = d^2 \times \zeta(2) \approx 41.12}
\end{equation}
derived from (i) $d = 5$ (the unique dimension, proven from atomic number theory), (ii) $\zeta(2) = \pi^2/6$ (Euler's Basel sum, the total lattice illuminance per channel).

\begin{remark}[A constant of the plateau, not the universe]
Even $25\pi^2/6$ is not a universal constant. It is a geometric identity of the rank-5 plateau --- the long-lived intermediate state that survives swap annihilation of repeated atomic dimensions. The universe itself is rank $N$ (arbitrarily large); the rank-5 plateau is a transient structure, albeit one that persists long enough for chemistry, biology, and observers. In the block universe, regions of rank $\gg 5$ (Big Bang) and rank $\to 1$ (heat death) coexist with the plateau. The ``constant'' $25\pi^2/6$ holds wherever and whenever the rank-5 structure is active --- which, for all practical purposes in our epoch, is everywhere we can observe.

The hierarchy of constants:

\begin{center}
\begin{tabular}{lcl}
\hline
Quantity & Status & Why \\
\hline
(none) & No universal constant & Universe $=$ rank-$N$ matrix \\
$25\pi^2/6$ & Plateau constant & Geometric identity of rank-5 \\
$1/\alpha_3, 1/\alpha_2, 1/\alpha_1$ & Full topological values & $\det(G_h)$ redistribution \\
$1/\alpha_\text{em} \approx 137$ & Local composite & $= \tfrac{5}{3}(1/\alpha_1) + (1/\alpha_2)$ \\
$1/137.036$ & Our address & $\varepsilon_0 \approx 0.0038$ here, now \\
\hline
\end{tabular}
\end{center}
\end{remark}

%=====================================================================
\subsection{Conclusion}
%=====================================================================

The Webb dipole is not a systematic error. It is not evidence for new fields or extra dimensions. It is the spatial non-uniformity of $\det(G_h)$ redistribution in $\CC^5$, visible because $\alpha_\text{em}$ is a composite of individual sector couplings whose geometric weights vary spatially.

The 25-year debate between ``$\alpha$ varies'' and ``$\alpha$ is constant'' is resolved: both are right, about different quantities. $\alpha_\text{em}$ varies. $\agut$ is constant. The difference is the $\det(G_h)$ redistribution among sectors.

The theory makes an immediately testable prediction: $\alpha_s$ must anti-correlate with $\alpha_\text{em}$ spatially, preserving $\sum\Delta_i = 0$. This can be checked with existing quasar data. A positive result would confirm the redistribution mechanism; a negative result would falsify it.

\begin{center}
\emph{The fine structure constant is not a constant.}

\emph{It is a $\det(G_h)$ shadow, cast differently in different places.}

\emph{$25\pi^2/6$ is the constant of our plateau --- the rank-5 era.}

\emph{The universe itself has no constants. Only structure.}
\end{center}
